\section{Biological Plausibility: ¿puede el cerebelo implementar lectura y escritura al estilo SDM?}

Tras la correspondencia teórica, la primera conferencia lleva la pregunta al nivel de la implementación: si la attention-like retrieval solo requiere activaciones de alta dimensión, un canal de escritura independiente y una pooled readout, ¿ya cuenta el cerebro con circuitos similares? El curso elige el cerebellum, no para afirmar que el cerebelo «es un Transformer», sino para buscar, elemento por elemento, componentes candidatos para address activation, value writing, superposition y majority-like pooling.

\subsection{Circuito abstracto: canal de direcciones y canal de escritura separados}

La diapositiva de Circuitry parte de una abstract neuron. El pattern de entrada determina qué neuronas de dirección hacen firing; otra vía, la value signal, determina qué deben almacenar esas neuronas; en la lectura, la query vuelve a activar las neuronas, cada una emite su value superpuesto y, al final, una pooled unit los agrega. Esta estructura hetero-associative corresponde directamente a la separación key/value de la attention.

\lecturefigure{slide-36-circuitry-sdm.jpg}{Circuito neuronal abstracto que implementa SDM}{Diapositiva V036 recuperada de la grabación oficial; 00:26:38}

\lecturefigure{slide-37-circuitry-mapping-final.jpg}{Estado final de la correspondencia del circuito abstracto con los cerebellar components}{Diapositiva V037 recuperada de la grabación oficial; 00:28:22}

\begin{knowledgebox}{Términos clave: papeles candidatos en el cerebellar circuit}
Las \term{mossy fibers} transportan la entrada; las \term{granule cells} forman una expansión dispersa de alta dimensión; las \term{parallel fibers} envían ampliamente la granule activity a las Purkinje cells; las \term{climbing fibers} aportan una señal de enseñanza/escritura intensa y relativamente independiente; las \term{Purkinje cells} agregan una gran cantidad de entradas e influyen en los deep cerebellar nuclei. El curso alinea estos papeles en el nivel funcional.
\end{knowledgebox}

\subsection{Anatomía del cerebelo y correspondencia, operación por operación, con SDM}

La diapositiva de anatomía muestra los microcircuitos repetidos y relativamente uniformes del cerebelo; la diapositiva final de mapping superpone las vías abstractas de read/write sobre cell types concretos. Las granule cells son enormemente numerosas y tienen conexiones dispersas, por lo que pueden funcionar como hard locations; el mossy input determina la activación; la señal intensa de la climbing fiber puede modificar una Purkinje synapse específica; la Purkinje cell recibe una cantidad enorme de parallel-fiber input y proporciona la pooled output.

\lecturefigure{slide-38-cerebellum-anatomy.jpg}{Tipos de neuronas y estructura de conexiones del cerebellum}{Diapositiva V038 recuperada de la grabación oficial; 00:31:04}

\lecturefigure{slide-39-sdm-cerebellum-mapping.jpg}{Comparación completa entre read/write de SDM y la cerebellar circuitry}{Diapositiva V039 recuperada de la grabación oficial; 00:34:00}

\teachervoice{El ponente afirma que alrededor del 70\% de las neuronas del cerebro se ubican en el cerebellum y advierte que posiblemente no solo participa en el control motor fino; sin embargo, también describe la evidencia sobre cognición de alto nivel como un área de investigación activa. Estas notas conservan ese tono exploratorio y no convierten el hecho de que una región «se ilumine» en la fMRI en un módulo de attention lingüística ya establecido.}

\subsection{Cómo pasar de «implementable» a «comprobado»}

La sesión de preguntas del curso marca el límite científico más importante: aunque el circuito pueda implementar SDM, eso no significa que el cerebro realmente lo use. Una verificación auténtica requiere registrar en tiempo real la sparse neuron activation, el synaptic update y el behavior change durante una associative task; idealmente, también escribir nuevas asociaciones y observar cómo se desvanecen las anteriores. El ponente propone la calcium/voltage imaging de alta velocidad como vía candidata; se trata de un diseño experimental, no de una conclusión ya obtenida.

\begin{warningbox}{Tres brechas de evidencia que no pueden omitirse}
Primero, que la conectividad anatómica permita cierto cómputo no significa que la dinámica lo ejecute; segundo, que la lectura y escritura local se parezcan no significa que todo el cerebro use el training objective de un Transformer; tercero, la analogía entre microzone y multi-head es atractiva, pero el ponente afirma explícitamente que no conviene llevarla demasiado lejos.
\end{warningbox}

En la Q\&A también se distingue entre la fast neural activity y la long-term synaptic plasticity, y se discute la compensación entre capacidad y fidelity que implica el radius: con un radio de activación grande, cada lectura agrega más neuronas y la estimación es más suave, pero aumenta la interferencia entre patterns; con un radio pequeño, la capacidad de discriminación es alta, pero una query con ruido puede activar muy pocas posiciones. En la attention, el parámetro correspondiente no es un radio físico fijo, sino la effective sharpness que producen conjuntamente las key/query norms y beta.

\lecturefigure{slide-40-first-talk-thank-you.jpg}{Resumen de la primera conferencia, referencias y agradecimientos}{Diapositiva V040 recuperada de la grabación oficial; 00:35:54}

\begin{importantbox}{Las conclusiones más sólidas y más débiles de la primera conferencia}
La conclusión más sólida es que, bajo condiciones explícitas, existe una aproximación derivable y verificable numéricamente entre la SDM intersection read y la softmax attention. Una conclusión más débil, aunque sugerente, es que el cerebellar circuit cuenta con componentes capaces de implementar operaciones similares de address/value/readout. Ambas deben presentarse por separado: la primera es evidencia matemática; la segunda, una hipótesis de neurociencia.
\end{importantbox}

\subsection{Resumen de la sección}

El mapping del cerebelo muestra una posible implementación biológica de una attention-like associative memory: la expansión dispersa, la señal de enseñanza independiente, el almacenamiento distribuido y la pooled readout tienen correspondencias funcionales. Sin embargo, la teacher voice más importante de la clase es que «implementable no equivale a verificado»: los experimentos deben seguir directamente los cambios de conectividad durante el aprendizaje y sus resultados conductuales.
